% Exact two-sided certificates for the doubly tilted CHSH family
% Every numbered statement is traced to a verifier claim in Appendix C (claim IDs from projects/bell/paper_belege.json).
% Numbers in the statements are the verifier's outputs, rounded outward where they bound Q.
\documentclass[10pt,twocolumn,a4paper]{article}
\usepackage[T1]{fontenc}
\usepackage[utf8]{inputenc}
\usepackage{lmodern}
\usepackage[english]{babel}
\usepackage{amsmath,amssymb,amsthm}
\usepackage[a4paper,top=22mm,bottom=24mm,left=17mm,right=17mm,columnsep=7mm]{geometry}
\usepackage{booktabs}
\usepackage{microtype}
\usepackage{fancyhdr}
\usepackage[hidelinks]{hyperref}
\usepackage{xcolor}

\pagestyle{fancy}
\fancyhf{}
\fancyhead[L]{\small Preprint (2026)}
\fancyhead[R]{\small Team Ninja Turtles: Exact certificates for the doubly tilted CHSH family}
\fancyfoot[C]{\small\thepage}
\renewcommand{\headrulewidth}{0pt}
\fancypagestyle{first}{\fancyhf{}\fancyfoot[C]{\small\thepage}}

\newtheorem{theorem}{Theorem}
\newtheorem{proposition}[theorem]{Proposition}
\newtheorem{lemma}[theorem]{Lemma}
\theoremstyle{definition}
\newtheorem{observation}{Numerical observation}
\theoremstyle{remark}
\newtheorem{remark}{Remark}

\newcommand{\Q}{\mathrm{Q}}
\newcommand{\Lc}{\mathrm{L}}
\newcommand{\NPA}[1]{\mathrm{NPA}_{#1}}
\newcommand{\ev}[1]{\langle #1\rangle}

\begin{document}
\thispagestyle{first}

\twocolumn[{%
\noindent{\small Preprint}\\[-1pt]
{\small Quantum Physics}\\[14pt]
\begin{center}
{\LARGE\bfseries Exact two-sided certificates for maximal quantum\\[3pt] violations of the doubly tilted CHSH inequality}\\[12pt]
{\large Team Ninja Turtles}\\[3pt]
{\small Hack-Nation 2026, Challenge 3 (Agentic Scientific Discovery)}\\[8pt]
{\small Preprint, October 2026}
\end{center}
\vspace{14pt}
}]

\noindent\textbf{Abstract}\enspace
We study the maximal quantum value $\Q(p)$ of the doubly tilted CHSH expression
$I_p = p\ev{A_1}+p\ev{B_1}+\ev{A_1B_1}+\ev{A_1B_2}+\ev{A_2B_1}-\ev{A_2B_2}$, whose classical bound is $\Lc(p)=2+2p$.
Lower bounds come from explicit quantum strategies, upper bounds from the Navascu\'es--Pironio--Ac\'in (NPA) hierarchy; both are
certified in exact rational arithmetic, the former by rational $\pm1$-observables and a rational state, the latter by a rational
dual certificate whose positivity is verified by an exact $LDL^{\mathsf T}$ factorisation. We prove that
$\Q(3/10)\in[2.938288331758,\,2.938288368082]$, $\Q(1/2)\in[3.136937311309,\,3.136937317689]$ and
$\Q(131/200)\in[3.358259513883,\,3.358259555321]$, each attained to within $10^{-7}$ by a qubit strategy. At $p=131/200$ the
optimum of the NPA level $1{+}AB$ exceeds $\Q$ by at least $3.25\cdot10^{-6}$, while level~2 is tight to $4.2\cdot10^{-8}$; so
level $1{+}AB$ is strictly not tight for this tilted CHSH expression. At $p=0.95$ and $p=0.99$ explicit strategies exceed $\Lc(p)$
exactly. Numerically, the verifier detects no gap of level $1{+}AB$ up to $p=0.62$, and the detected gap grows monotonically from $p=0.652$ on.
The questions were posed, the experiments planned and the claims formulated by language-model agents; only claims that passed the
exact verifier are reported, and the failures are listed. No closed form for $\Q(p)$ was found.

\section{Introduction}

For a bipartite Bell expression the classical bound is a finite maximisation, whereas the maximal quantum value is an optimisation
over states and measurements in Hilbert spaces of unbounded dimension. For the CHSH expression~\cite{chsh} the latter equals
$2\sqrt2$~\cite{tsirelson}, and for many other expressions it is known only numerically. The standard tools bound it from two
sides: explicit strategies give lower bounds, and the NPA hierarchy of semidefinite programs gives a decreasing sequence of upper
bounds~\cite{npa07,npa08}, which in the limit yields the quantum value~\cite{avis}. Whether a given level is already tight
depends on the expression; for the inequality $I_{3322}$~\cite{cg} finite-dimensional strategies are conjectured not to attain the
maximum at all~\cite{pv}.

Tilting CHSH by marginal terms changes this picture. For one-sided tilts $\beta\ev{A_1}+\text{CHSH}$ the quantum value
$\sqrt{8+2\beta^2}$ is known~\cite{amp}, and operations that change the classical bound while leaving the quantum
bound invariant have been studied systematically~\cite{epping}. In this note we consider the symmetric two-sided tilt
\begin{equation}\label{eq:Ip}
I_p = p\ev{A_1} + p\ev{B_1} + \ev{A_1B_1}+\ev{A_1B_2}+\ev{A_2B_1}-\ev{A_2B_2},
\end{equation}
with $\pm1$-valued observables $A_x,B_y$ and $0\le p\le 1$, and we ask three questions: what is $\Q(p)$, do qubits suffice, and from
which $p$ on is the NPA level $1{+}AB$ no longer tight?

Our emphasis is on certification. Every value we state as a theorem is enclosed between a lower bound that is the exact Bell value
of an explicit rational strategy and an upper bound that is the exact value of a rational dual certificate. Floating-point
computations only propose candidates. The study was carried out by an agentic laboratory in which language-model agents proposed
questions, experiments and claims, and a fixed program decided what is true (Section~\ref{sec:lab}). We report the negative
outcomes alongside the positive ones.

Two remarks delimit the content. First, the theorems concern finitely many parameter values; no closed form for $\Q(p)$ is
claimed. Second, our literature search was automated and restricted to abstracts; it did not return a treatment of~\eqref{eq:Ip},
but we make no claim of novelty.

\section{Setting and certificates}\label{sec:setting}

A two-outcome Bell expression in correlator form is
$I = c_0 + \sum_x a_x\ev{A_x} + \sum_y b_y\ev{B_y} + \sum_{x,y} c_{xy}\ev{A_xB_y}$ with rational coefficients. For~\eqref{eq:Ip},
$a=b=(p,0)$, $c=\left(\begin{smallmatrix}1&1\\1&-1\end{smallmatrix}\right)$ and $c_0=0$.

\begin{lemma}\label{lem:L}
For $p\ge 0$ the classical bound of~\eqref{eq:Ip} is $\Lc(p)=2+2p$.
\end{lemma}
\begin{proof}
For deterministic values $a_x,b_y\in\{\pm1\}$, if $b_1=b_2=b$ the expression equals $p(a_1+b)+2a_1b\le 2p+2$; if $b_1=-b_2$ it
equals $p(a_1+b_1)+2a_2b_1\le 2p+2$. All values $+1$ attain $2+2p$.
\end{proof}

The verifier certifies three kinds of statements; the details are in Appendix~\ref{app:cert}.

\emph{Classical bound.} Exact enumeration of all deterministic strategies in rational arithmetic.

\emph{Lower bound on $\Q$.} A real strategy of local dimension $d$ is found by see-saw iteration. The verifier rounds the
eigenvectors of each observable with eigenvalue $-1$ to a rational matrix $V$ and sets $A=I-2V(V^{\mathsf T}V)^{-1}V^{\mathsf T}$,
which is exactly a symmetric involution; the state is rounded to a rational vector and normalised exactly. The Bell value of the
resulting strategy is then a rational number and a valid lower bound on $\Q$.

\emph{Upper bound on $\Q$.} For an NPA level with moment matrix indexed by a word set $S$, any symmetric $Z\succeq0$ with
$\langle Z,F_w\rangle=-c_w$ for all words $w\ne 1$ gives $\Q\le c_0+\langle Z,F_1\rangle$. The verifier rationalises the
numerical dual solution, restores the linear constraints exactly and certifies $Z+\varepsilon I\succ0$ by an exact
$LDL^{\mathsf T}$ factorisation; since the diagonal belongs to the identity class, the bound becomes
$c_0+\langle Z,F_1\rangle+\varepsilon|S|$. We use the levels $1{+}AB$, $2$ and $2{+}AAB$ (Appendix~\ref{app:cert}).

We call $\Q$ \emph{proved} at a parameter value when the two bounds differ by at most $10^{-6}$. To show that a level is
\emph{not} tight we also need a lower bound on the optimum of the level itself; it is obtained from a rational feasible moment
matrix (Appendix~\ref{app:cert}).

Before the study, the verifier passed a self-test of 22 cases, 10 true statements and 12 false statements or rule violations.
The true anchors include the CHSH value $2.828427125$, the classical bound $6$ of the chained inequality with four
settings~\cite{bc}, the classical bound $0$ and an NPA upper bound $0.2509$ for $I_{3322}$, and the formula $\sqrt{8+2p^2}$
of~\cite{amp} at five points. The false statements sit just beyond known values, e.g.\ an upper bound $2.8284270$ for CHSH, and
the rule violations include a tolerance supplied with the claim, which the verifier ignores.

\section{Results}\label{sec:results}

\begin{theorem}[computer-assisted]\label{thm:values}
The maximal quantum value of~\eqref{eq:Ip} satisfies
\begin{align*}
2.938288331758 &\le \Q(3/10) \le 2.938288368082,\\
3.136937311309 &\le \Q(1/2) \le 3.136937317689,\\
3.358259513883 &\le \Q(131/200) \le 3.358259555321.
\end{align*}
In each case the lower bound is attained by a strategy with local dimension $d=2$.
\end{theorem}

The upper bounds come from the NPA levels $1{+}AB$, $2{+}AAB$ and $2$, respectively. The widths are $3.6\cdot10^{-8}$,
$6.4\cdot10^{-9}$ and $4.2\cdot10^{-8}$, so in particular $\Q$ is proved at these points in the sense of
Section~\ref{sec:setting}, and qubit strategies attain the quantum value to within these widths. Independent see-saw searches of
the verifier in dimensions up to $4$ found no larger values.

\begin{theorem}[computer-assisted]\label{thm:npa}
At $p=131/200$ the optimum of the NPA level $1{+}AB$ for~\eqref{eq:Ip} is at least $3.3582628097$, whereas
$\Q(131/200)\le 3.358259555321$. Hence level $1{+}AB$ exceeds the quantum value by at least $3.25\cdot10^{-6}$, while level~$2$
is tight to within $4.2\cdot10^{-8}$.
\end{theorem}

The first number is the value of a rational positive definite moment matrix that satisfies all constraints of level $1{+}AB$
exactly; the second is the upper bound of Theorem~\ref{thm:values}, certified at level~2.

\begin{proposition}[computer-assisted]\label{prop:viol}
At $p=0.99$ an explicit qubit strategy attains $I_p = 3.98000132884\ldots > 3.98 = \Lc(0.99)$, and
$\Q(0.99)\le 3.98000623487$. At $p=0.95$,
$3.90016389936 \le \Q(0.95) \le 3.90019981473$, so $\Q(0.95)>\Lc(0.95)=3.9$.
\end{proposition}

The violations at these points are small, $1.3\cdot10^{-6}$ and $1.6\cdot10^{-4}$, but exact. At $p=0.95$ the bounds differ by
$3.6\cdot10^{-5}$, so $\Q(0.95)$ is bounded but not proved. At $p=1$ the verifier's best strategy attains exactly $4=\Lc(1)$ and no
violation was found; that $\Q(1)=4$ is not shown.

\begin{observation}\label{obs:gap}
Let $g(p)$ denote the verifier's lower bound for the optimum of level $1{+}AB$ minus the best exactly certified strategy value.
\begin{center}
\begin{tabular}{@{}lcccccc@{}}
\toprule
$p$ & $0.5$ & $0.62$ & $0.652$ & $0.6532$ & $0.6533$ & $0.655$\\
\midrule
$g(p)$ & $-1.7\text{e-}8$ & $-2.3\text{e-}8$ & $1.8\text{e-}7$ & $1.03\text{e-}6$ & $1.12\text{e-}6$ & $3.30\text{e-}6$\\
\bottomrule
\end{tabular}
\end{center}
Up to $p=0.62$, $|g|$ is below $10^{-7}$, the accuracy of the SDP solvers; from $p=0.652$ on $g$ is positive and grows
monotonically over the points examined.
\end{observation}

A threshold defined by a fixed tolerance, such as ``$g>10^{-6}$'', therefore depends on that tolerance; we do not locate the onset
of non-tightness. Only at $p=131/200$ does $g$ compare with a proved quantum value (Theorem~\ref{thm:npa}). At $p=0.72$ level~2
shows no detectable gap: the corresponding bound for level~2 is $-5.7\cdot10^{-8}$.

\begin{remark}[status of the computer-assisted parts]\label{rem:status}
All statements labelled computer-assisted rest on exact rational arithmetic (Python \texttt{fractions}). Floating-point
see-saw iterations and the SDP solvers (Clarabel, SCS via CVXPY) only produce candidates; they are not part of the trusted base.
The trusted base consists of the enumeration, the exact evaluation of the rounded strategy, the exact restoration of the dual
constraints and the exact $LDL^{\mathsf T}$ test. The bounds in Theorems~\ref{thm:values} and~\ref{thm:npa} are the verifier's
exact values rounded outward to the digits shown.
\end{remark}

\section{The agentic laboratory}\label{sec:lab}

The study was run by the verifier-gated laboratory of~\cite{repo}. A literature scout queried arXiv and Europe PMC (320 sources,
84 extracted findings, 83 of them with a quotation found verbatim in the abstract by code). An integrator chose one question per
round by expected information gain and preregistered success and stopping criteria before any experiment. A cascade of researcher
agents (two smaller, two larger language models) planned experiments, read their results and stated typed claims; the cascade
stopped at the first claim that passed the verifier. A red team then tried to refute each passed claim with up to two
counter-claims of the same types, and a learning agent proposed follow-up questions. Language models never decide what is true; a
claim stands only if the verifier passes it and no counter-claim does.

The laboratory ran eight rounds at a cost of $4.31$\,USD. Seven claims passed the verifier; one of them was contested by the red
team (Section~\ref{sec:neg}), and one round produced no claim that passed. The manuscript was drafted by a writing agent under a
gate that rejects every sentence whose numbers do not occur in a cited verified claim, and was then edited for presentation. All
numbered statements are traced to their claims in Appendix~\ref{app:prov}.

\section{Negative results and preregistration}\label{sec:neg}

\emph{Contested claim.} In the last round the agents claimed that the transition of level $1{+}AB$ lies between $p=0.6532$ and
$p=0.6533$, defined by $g>10^{-6}$. The red team showed $g(0.6532)\ge1.025\cdot10^{-6}$, refuting the lower end. This is what led
to Observation~\ref{obs:gap}: the location depends on the tolerance.

\emph{Failed round.} A coarse scan for the same transition over $p\in\{0.6,\dots,0.99\}$ produced no claim that passed.

\emph{Formula.} Two rounds searched for a closed form of $\Q(p)$, including ansätze of the form $\Q^2=8+c_1p^2+c_2p^4$. The
agents submitted fitted candidates such as $\sqrt{8+7p^2+\tfrac32p^4}$; the verifier rejected it at its first support point, where
it predicts $2.840800943$ against the certified enclosure $\Q(1/10)\in[2.840515577,\,2.840515578]$. No formula passed.

\emph{Preregistration.} Before the run we registered two hypotheses. H6a, that at least two of three anchor questions (CHSH,
chained, one-sided tilt) would be confirmed, was not met: the integrator selected none of them. H6b, that at least one certified
statement concerns an expression outside the named families, was met with six confirmed claims on~\eqref{eq:Ip}.

\section{Discussion}

The doubly tilted CHSH expression is a simple example in which the cheapest useful NPA level stops being tight while the next one
remains tight, and in which this can be decided with exact certificates on both sides. The certificates are simple: a rational
strategy for the lower bound and a rational positive semidefinite dual matrix for the upper bound; both can be rechecked
independently of the solvers that produced them. The results are consistent with the literature as far as we checked it: the
one-sided formula of~\cite{amp} is reproduced by the same verifier, and for the simplest scenario a coincidence of levels $1{+}AB$
and $2$ has been reported for correlations under a plausible extremality criterion~\cite{ishizaka}. Our Theorem~\ref{thm:npa}
concerns an expression with marginal terms; whether the scope of that statement covers such expressions we did not check beyond
the abstract, and we do not claim a contradiction.

Several questions remain open. (i) A closed form for $\Q(p)$, or at least an algebraic equation satisfied by it, is missing; the
certified values of Theorem~\ref{thm:values} can test any candidate. (ii) The onset of non-tightness of level $1{+}AB$ is not
located; in Observation~\ref{obs:gap} the detected gap leaves the noise level between $p=0.62$ and $p=0.652$, but a proof of
the onset requires certified values of $\Q$ and of the level on both sides of it. (iii) Whether $\Q(p)>\Lc(p)$ for every $p<1$ is open; Proposition~\ref{prop:viol} shows it at two points close to
$1$, where the violation is of order $10^{-6}$ to $10^{-4}$. (iv) Whether qubits are optimal for all $p$ is open; they are
optimal to within the stated widths at the three certified points.

\appendix
\section{The certificates in detail}\label{app:cert}

\emph{Words and moment matrices.} Letters are $A_x,B_y$ with $A_x^2=B_y^2=1$ and $[A_x,B_y]=0$. Words are reduced by moving all
$A$-letters before all $B$-letters and cancelling equal neighbours. For real strategies $\ev{w}=\ev{w^{\mathsf T}}$, so a word
is identified with its reverse. Restricting to real symmetric moment matrices loses nothing, since the real part of a feasible
complex moment matrix is feasible with the same objective value. The word sets are $1{+}AB=\{1,A_x,B_y,A_xB_y\}$,
$2=1{+}AB\cup\{A_xA_{x'},B_yB_{y'}\}$ and $2{+}AAB=2\cup\{A_xA_{x'}B_y,A_xB_yB_{y'}\}$.

\emph{Dual certificate.} The primal is $\max\sum_w c_w y_w$ subject to $\Gamma(y)=F_1+\sum_{w\ne1}y_wF_w\succeq0$. For
$Z\succeq0$ with $\langle Z,F_w\rangle=-c_w$, $0\le\langle Z,\Gamma\rangle=\langle Z,F_1\rangle-\sum_{w\ne1}c_wy_w$. The numerical
$Z$ is rounded to denominators $10^{10}$; each class constraint is restored exactly by spreading the defect uniformly over the
entries of the class; $\varepsilon$ is the smallest value in a fixed list above $\max(0,-\lambda_{\min})$ for which the exact
$LDL^{\mathsf T}$ factorisation of $Z+\varepsilon I$ has positive pivots.

\emph{Feasible moment matrix.} For a lower bound on the optimum of a level, the numerical primal $\Gamma$ is averaged over each
class, rounded, its identity class set to $1$, and replaced by $\Gamma_t=(1-t)\Gamma+tI$ with the smallest $t$ in a fixed list
for which $\Gamma_t\succ0$ exactly. Since $I$ is feasible, $\Gamma_t$ is feasible, and its objective is a lower bound.

\emph{Strategy.} The see-saw iteration alternates the top eigenvector of the Bell operator with $A_x=\operatorname{sign}(M_x)$ for
the effective operators $M_x$. The verifier's own search uses fixed seeds disjoint from those available to the agents.

\section{Self-test of the verifier}\label{app:self}

The verifier must pass its self-test before the laboratory starts. Of the 22 cases, the true ones are listed in
Section~\ref{sec:setting}. The false ones include the classical bound $201/100$ for CHSH, a lower bound $2.8284272$ and an upper
bound $2.8284270$ for CHSH (just beyond $2\sqrt2$), the value $0.25$ for $I_{3322}$ as a proved quantum value (the gap is not
closed), the formula $\sqrt{8+p^2}$, a strategy for another expression, an unknown strategy identifier and a formula with fewer
than three support points. During development the self-test also corrected an expectation of ours: the upper bound $0.2509$ for
$I_{3322}$, entered as false, is true and is certified at level $2{+}AAB$.

\section{Provenance}\label{app:prov}

Claim identifiers refer to \texttt{projects/bell/paper\_belege.json} in~\cite{repo}.

\begin{center}\small
\begin{tabular}{@{}ll@{}}
\toprule
Statement & Claims\\
\midrule
Thm.~\ref{thm:values}, $p=3/10$ & C-bell-R4, -R4-RT2\\
Thm.~\ref{thm:values}, $p=1/2$ & C-bell-R1, -R3, -R3-RT2, -R1-RT1\\
Thm.~\ref{thm:values}, $p=131/200$ & C-bell-R6, -R6-RT2\\
Thm.~\ref{thm:npa} & C-kombi1 (from -R5, -R6)\\
Prop.~\ref{prop:viol} & C-bell-R2, -R2-RT1, -R2-RT2\\
$p=1$ & C-bell-R1-RT2\\
Obs.~\ref{obs:gap} & C-kombi2, -R5-RT1, -R5-RT2\\
Contested claim & C-bell-R8, -R8-RT1, -R8-RT2\\
Rejected formula & \texttt{runde4.json} (trace)\\
Failed round & C-neg1\\
Preregistration & C-H6a, C-H6b\\
Self-test & C-selbsttest\\
Laboratory statistics & C-methode\\
Literature & C-lit5, -6, -9, -10, -13, -14\\
\bottomrule
\end{tabular}
\end{center}

\section*{Declarations}

\noindent\textbf{Acknowledgements}\enspace Questions, experiments and claims were produced by language-model agents (Claude,
Anthropic) inside the laboratory described in Section~\ref{sec:lab}; the manuscript was drafted by such an agent and edited with
its assistance. Every mathematical statement was accepted only by the exact verifier. The team is responsible for the content.

\smallskip
\noindent\textbf{Code and data availability}\enspace All code, the verifier and its self-test, the stored strategies, the
laboratory state, the preregistration and the language-model transcripts are in~\cite{repo}, branch \texttt{paper-bell}.

\smallskip
\noindent\textbf{Competing interests}\enspace The authors declare no competing interests.

\begin{thebibliography}{99}\small
\bibitem{chsh} J.F. Clauser, M.A. Horne, A. Shimony, R.A. Holt, Proposed experiment to test local hidden-variable theories.
Phys. Rev. Lett. 23, 880--884 (1969). \url{https://doi.org/10.1103/PhysRevLett.23.880}
\bibitem{tsirelson} B.S. Cirel'son, Quantum generalizations of Bell's inequality. Lett. Math. Phys. 4, 93--100 (1980).
\url{https://doi.org/10.1007/BF00417500}
\bibitem{bc} S.L. Braunstein, C.M. Caves, Wringing out better Bell inequalities. Ann. Phys. 202, 22--56 (1990).
\url{https://doi.org/10.1016/0003-4916(90)90339-P}
\bibitem{cg} D. Collins, N. Gisin, A relevant two qubit Bell inequality inequivalent to the CHSH inequality. J. Phys. A 37,
1775--1787 (2004). \url{https://doi.org/10.1088/0305-4470/37/5/021}
\bibitem{npa07} M. Navascu\'es, S. Pironio, A. Ac\'in, Bounding the set of quantum correlations. Phys. Rev. Lett. 98, 010401
(2007). \url{https://doi.org/10.1103/PhysRevLett.98.010401}
\bibitem{npa08} M. Navascu\'es, S. Pironio, A. Ac\'in, A convergent hierarchy of semidefinite programs characterizing the set of
quantum correlations. New J. Phys. 10, 073013 (2008). \url{https://doi.org/10.1088/1367-2630/10/7/073013}
\bibitem{avis} D. Avis, S. Moriyama, M. Owari, From Bell inequalities to Tsirelson's theorem: a survey (2008).
\url{https://arxiv.org/abs/0812.4887}
\bibitem{pv} K.F. P\'al, T. V\'ertesi, Maximal violation of a bipartite three-setting, two-outcome Bell inequality using
infinite-dimensional quantum systems. Phys. Rev. A 82, 022116 (2010). \url{https://doi.org/10.1103/PhysRevA.82.022116}
\bibitem{amp} A. Ac\'in, S. Massar, S. Pironio, Randomness versus nonlocality and entanglement. Phys. Rev. Lett. 108, 100402
(2012). \url{https://doi.org/10.1103/PhysRevLett.108.100402}
\bibitem{epping} M. Epping, H. Kampermann, D. Bru\ss, Optimisation of Bell inequalities with invariant Tsirelson bound (2014).
\url{https://arxiv.org/abs/1402.0378}
\bibitem{ishizaka} S. Ishizaka, NPA hierarchy and extremal criterion in the simplest Bell scenario. Entropy 27, 182 (2025).
\url{https://doi.org/10.3390/e27020182}
\bibitem{repo} Team Ninja Turtles, Verifier-gated discovery lab, source code and data.
\url{https://github.com/NinjaTurtlesHackathons/Hacknation} (branch \texttt{paper-bell}).
\end{thebibliography}

\end{document}
